\section{The base reduction with varying depth inputs}
\label{sec:base}

We prove Theorem~\ref{thm:depth-local} by adapting the local-form
construction of \cite[Sections~3--10]{H}.  The scalar analytic and
finite combinatorial estimates remain those of H.  The row estimates
of Section~\ref{sec:rows} replace the uses of one original input at
different depths.  This section constructs the tuple to which the
later localization applies: cyclic approximations first give a tuple
\(X\), delayed high fits replace it by \(Y\), and the decomposition
\(Y=B+U\) leaves three active patterns.

Fix the root \(J\), finite tree \(\mathcal T\), original summation
collection \(\mathcal I\), and four input schedules of
Theorem~\ref{thm:depth-local}.  Write
\(s(I)=\log_2(|J|/|I|)\).  On every subroot \(R\), the spatial origin
\(s_R\) of Section~\ref{sec:rows} is \(s(R)\); a restart does not reset
depth.  The actual augmented input is \(F_{j,s}=(f_{j,s},0)\).
Finitely needed temporal indices beyond the schedule use its nearest
endpoint value, while the spatial inputs remain zero outside \(J\).
In the local forms put \(y_j=x+jt\), \(0\le j\le3\), and use the
normalized norms of Definition~\ref{def:transfer-local-kernels}.
Only physical components enter a form.  The notation \(\Heq{n}\)
denotes Equation~\((n)\) of H.

Fix provisionally one \(q\) with \(2<q<3\).  Section~\ref{sec:completion}
will choose a single value satisfying both arguments' degree restrictions
and run H's fixed-input proof at that exponent.  Polynomial degrees below are
absolute; their coefficients and the lower accuracy threshold may
depend on \(q,c_0,C_0\).  Accuracy indices \(t\) are independent of
spatial depths \(s,h\).

\subsection{Inputs, source columns, and the row estimate}

Lemma~\ref{lem:depth-input-size} and Equation~\eqref{eq:row-variation}
give, on every subroot \(R\subseteq J\),
\begin{gather}
 \|F_{j,s}\|_{p,I}\le1,\qquad
 \|\mathbf1_R F_{j,s}\|_2\le |R|^{1/2}
 \quad(1\le p\le\infty),                                      \label{eq:h-input-size}\\
 \sum_{[r,t]}\|\mathbf1_R(F_{j,t}-F_{j,r})\|_2^2
       \le5a|R|.                                               \label{eq:h-input-variation}
\end{gather}
The first bound holds on every descendant \(I\) and every prescribed
or finitely extended input index.  The second holds for depth intervals
with edge overlap at most \(a\).  Restriction to a subroot, a
nondecreasing subsequence, and repeated reads preserve these
conclusions.  Each slot retains its own full block schedule.  We also
use the \emph{Local absolute estimate} Lemma of \cite[Section~2]{H}:
\[
 |H_I(z_0,z_1,z_2,z_3)|
 \le C\prod_{j=0}^3\|z_j\|_{p_j,I},
 \qquad p_j\ge1,\quad \sum_jp_j^{-1}\le2.
\]

The scalar labels are restrictions to their birth cells of the
bounded quadratic chart atoms in the \emph{Quadratic chart atom}
Definition of \cite[Section~4]{H}.  We use exactly that
source-defined class and complexity: \(H\ge e^2\) bounds the chart
dimension by \(\log H\), the rational matrix height and amplitude
support by \(H\), and the amplitude derivatives by the fixed
bounds in that Definition.  The real quadratic, linear, and
horizontal oscillation parameters are unrestricted.  A bounded
atom has supremum norm at most one, with scalar normalization
factors kept in the external coefficients.  The
\emph{Elementary chart operations} Lemma of \cite[Section~4]{H}
permits conjugation, bounded products, and affine changes of the
real argument with a quasipolynomial complexity change.  We write
\(e(v)=\exp(2\pi i v)\).

For a birth label \(v\), write
\(\Phi_v=(\varphi_v,\sqrt{\epsilon_v}\mathbf e_v)\), with the
persistent birth cell and distinct private direction of
Section~\ref{sec:rows}.  A joint budget \(D\ge2\) means exactly
Equation~\eqref{eq:row-catalog-budget}, including every available
label on a path.  Lemma~\ref{lem:row-catalog}, equivalently the
\emph{Coefficient and projection bounds} Lemma of
\cite[Section~3]{H}, gives for constant coefficients on a cell \(I\)
where the labels are available
\begin{equation}
 D^{-1}\|c\|_{\ell^2}\le
 \left\|\sum_v c_v\Phi_v(y)\right\|
 \le D^{3/2}\|c\|_{\ell^2}
 \quad\text{for almost every }y\in I.                         \label{eq:h-column-bounds}
\end{equation}
Let \(P_{I,\mathcal V}\) denote the local orthogonal projection onto
the constant-coefficient span of an available mask \(\mathcal V\).
The same source Lemma gives
\[
 \|P_{I,\mathcal V}z\|_{\infty,I}\le D^5\fint_I\|z\|,
 \qquad
 \fint_I\bigl(f-(P_{I,\mathcal V}(f,0))_{\rm phys}\bigr)
       \overline{\varphi_v}=\epsilon_v c_v,
\]
where \(c_v\) is that projection's coefficient.  Inherited masks
retain their admitted labels on children; their full-partition direct
sums are nested orthogonal projections.  At a summation cut, the
continuation convention of \cite[Section~3]{H} retains the old labels
and locally refits them on every finer cell needed for the old
summands.  A new attempt may use a separate family.  On a subroot,
Lemma~\ref{lem:row-catalog} rebases inherited labels without changing
their private constants.  These are local refits of the retained
spans.

For every fixed realized finite row array and catalog satisfying
Lemma~\ref{lem:row-depth}, with ordinary Bessel norm \(B\), that Lemma
gives
\begin{equation}
 \sum_\nu\|B_\nu F_{j,s(\nu)}\|^2
 \le C B^2\log^2(2D)|R|.                                     \label{eq:h-row-use}
\end{equation}
The read \(s(\nu)\) is common to its round and nondecreasing between
rounds.  The realized rows and catalog may have been selected using
the schedules; they are fixed while the ordinary operator inequality
is tested on arbitrary inputs.

\subsection{One capped trial and cyclic adaptive fits}

For a physical residual \(r\) in slot \(j\), let
\(\mathfrak d_{j,I}(r)\) be the supremum of the form against two
normalized \(L^2(I)\) inputs and one normalized \(L^q(I)\) input in
the other slots, with the \(L^q\) slot arbitrary.  The
\emph{Local detection} Lemma of \cite[Section~4]{H} says that if
\(C_1\ge1\), \(0<\delta<1/2\), \(\|r\|_{2,I}\le C_1\), and
\(\mathfrak d_{j,I}(r)>\delta\), there is a bounded chart atom
\(\varphi\) such that
\[
 \left|\fint_I r\overline\varphi\right|\ge\rho,\qquad
 \log(1/\rho),\ \log(\operatorname{complexity}(\varphi))
 \le C_q\bigl(1+\log C_1+\log(1/\delta)\bigr)^C.
\]
We may decrease \(\rho\) to lie in \((0,1]\).  The form norm of two
specified factors below is the supremum against two normalized
\(L^2(I)\) inputs in their complementary slots.

\begin{lemma}[One capped insertion trial]
\label{lem:hbase-trial}
Let \(R\) be an attempt root with a restricted normalized input
schedule \(F_s\).  Let the candidate request cells belong to a fixed finite
original request tree below \(R\).  Fix
\(0<\rho\le1\), \(D_0\ge2\), and an integer cap \(T\ge1\).
Process candidates in nondecreasing request-cell depth, allowing
repetitions at one cell.  Stop before the first request that would
exceed \(T\) successful requests on its spatial path, and remove
that cut cell and its candidate descendants from the trial.

For each successful request \(v\) at \(I_v\), suppose there is a
set \(\mathcal R_v\) of charged row indices, with these sets pairwise
disjoint, such that
\[
 \sum_{\nu\in\mathcal R_v}
       \|\Delta_\nu F_{s(\nu)}\|_2^2
 \ge c\rho^2|I_v|.
\]
Assume all charged rows are a Bessel-norm-one subfamily of one
realized finite nested sequence of full-partition projections and
satisfy every hypothesis of Lemma~\ref{lem:row-depth}.  Suppose its
completed catalog has budget
\[
 D_T\le C D_0(T+1)^C,
\]
where \(D_0\) includes all predetermined columns, including their
future admissions on this attempt.  Then the disjoint first-excess
cells \(\mathcal C_T\), which are original request-tree cells, obey
\begin{equation}
 \sum_{I\in\mathcal C_T}|I|
 \le C T^{-1}\rho^{-2}\log^2(2D_T)|R|.                        \label{eq:hbase-trial-cut}
\end{equation}
If
\[
 a=1+\log(2D_0)+\log(1/\rho),\qquad
 T=\left\lceil C_*\rho^{-4}(1+a)^8\right\rceil,
\]
then a sufficiently large fixed \(C_*\) makes the right side at
most \(|R|/4\), and \(\log(2T)\le C(a+\log C_*)\).
\end{lemma}

\begin{proof}
Freeze the completed rows and catalog.  By their disjoint ownership
and Equation~\eqref{eq:h-row-use},
\[
 c\rho^2\sum_v|I_v|
 \le\sum_{\nu\ {\rm charged}}\|\Delta_\nu F_{s(\nu)}\|_2^2
 \le C\log^2(2D_T)|R|.
\]
Set \(n(y)=\#\{v:y\in I_v\}\).  Earlier request cells meeting a
candidate contain it or equal it, because request depths are
nondecreasing.  Thus their count is constant on that candidate.
Each first-excess cell lies in \(\{n\ge T\}\), and removing its
subtree makes these cells disjoint.  Tonelli's identity
\(\int_R n=\sum_v|I_v|\) proves
Equation~\eqref{eq:hbase-trial-cut}.

For the stated cap, the budget bound gives
\(\log(2D_T)\le C'(a+\log C_*)\).  Consequently the coefficient
in Equation~\eqref{eq:hbase-trial-cut} is at most
\[
 C C_*^{-1}\rho^2(1+a)^{-8}(a+\log C_*)^2,
\]
which is at most \(1/4\) for fixed sufficiently large \(C_*\).
The logarithmic cap bound follows from the same formula.
\end{proof}

We realize the Lemma's nested sequence in both applications by one
rule.  At each projection depth, refine and locally refit the entire
root partition, admit the predetermined columns prescribed there,
and then perform simultaneous insertion rounds.  Cells with no
request retain their projections.  The next depth refines the entire
partition, including branches with no later tests for this attempt.
Refinements, admissions, and identity transitions remain in the
nested sequence.  Splitting an insertion increment on all cells of
its partition preserves its joint Bessel norm, by the splitting
identity in Corollary~\ref{cor:row-projections}.  The following two
applications specify the charged supports, observations, and reads.

Apply this construction first to the
\emph{Inherited adaptive approximation} Lemma of
\cite[Section~5]{H}, at accuracy \(e^{-k}\), for the fixed provisional
\(q\) and sufficiently large \(k\).  Every attempt starts with an
empty fit catalog.  On an assigned original candidate cell \(I\) of
depth \(s\), test
\[
 r=f_{j,s}-(PF_{j,s})_{\rm phys},
\]
where \(P\) is the current local projection.  Contraction gives
\(\|r\|_{2,I}\le2\).  On failure of the \(\mathfrak d_{j,I}\)
test, Local detection supplies a bounded atom with moment \(\rho\),
where \(\log(1/\rho)\le k^{O(1)}\).  Adjoin
\((\varphi,\mathbf e_v)\) with a fresh private direction and refit.
For the new projection \(P^+\),
\begin{equation}
 \|(P^+-P)F_{j,s}\|_{L^2(I)}^2
 \ge
 \frac{\left|\int_I
  (f_{j,s}-(PF_{j,s})_{\rm phys})\overline\varphi\right|^2}
      {\|(\varphi,\mathbf e_v)\|_{L^2(I)}^2}
 \ge c\rho^2|I|.                                               \label{eq:hbase-fresh-gain}
\end{equation}
Indeed, projection of the fresh column off the old span preserves
its pairing with the old residual and can only reduce its norm.

Here each request owns its insertion row on \(I\).  The support and
observation partitions are the full depth-\(s\) partition, and every
round at this depth reads \(F_{j,s}\).  Their refinement and read
order therefore satisfy Lemma~\ref{lem:row-depth}.  The initial
budget \(D_0\) is fixed, and a path cap \(T\) adds at most \(T\)
bounded columns of private weight one, giving the required polynomial
budget \(D_T\).  Lemma~\ref{lem:hbase-trial} makes the cut cells have
total length at most one quarter of the attempt root.

On every child root the restricted schedule is still normalized;
the fit again starts with no columns and has the same detection,
gain, full-partition chronology, and budget.  The one-trial Lemma
therefore applies anew.  Its length bound prevents a child from
cutting itself, so all new roots are proper original request-tree
descendants.  The finite recursion terminates, and the attempt-root
lengths sum to at most
\(\sum_{n\ge0}4^{-n}|R|=4|R|/3\), as in H's packing argument.
Each original summand belongs to one attempt.  The old projection
family continues by inheritance and local refitting only for its
old summands and their operator estimates.

On a retained cell the fit has passed:
\[
 g_I=P_IF_{j,s},\qquad
 \mathfrak d_{j,I}\bigl(f_{j,s}-(g_I)_{\rm phys}\bigr)\le e^{-k},
 \qquad \|g_I\|_{2,I}\le1.
\]
Within a segment, masks are inherited.  The cap and Local detection
bound path counts and atom complexities by \(\exp(k^{O(1)})\).
Private weights are one, so the coefficient \(\ell^2\) norm is at
most the normalized projection norm; Cauchy--Schwarz over the path
labels gives the same \(\exp(k^{O(1)})\) physical supremum bound.
All these statements hold anew on a normalized subroot.

We now use these fits cyclically in all four slots.  Let
\(\iota(t)=(t-1)\bmod4\) be the slot updated at accuracy level \(t\).
Choose an absolute \(A>1\), a sufficiently large initial \(k_1\),
and set \(k_{t+1}=k_t^A\).  Write \(g_{t,I}=P^t_IF_{\iota(t),s}\)
for the fit at level \(t\), and set \(g_{t-4,I}=0\) before that
slot's first update.  For any fixed bounded number of levels,
intersect their segment families.  Every common root is a root of
one constituent family or the original root, so their total length
remains bounded.

\begin{lemma}[The two earliest increments on depth inputs]
\label{lem:hbase-X-output}
On each original summation cell \(I\), put \(s=s(I)\), its original
spatial depth in \(J\).  At a finite terminal level \(n\), write \(g_{e,n,I}\) for the last
fit in slot \(e\) at or before \(n\), or zero if there is none, and
put \(r_{e,n,I}=F_{e,s}-g_{e,n,I}\).  For \(u<t\) with
\(\iota(u)=i\ne j=\iota(t)\), let \(a,b\) be the other slots and set
\begin{equation}
 \begin{aligned}
 X_i^{u,t}&=g_{u,I}-g_{u-4,I},&
 X_j^{u,t}&=g_{t,I}-g_{t-4,I},\\
 X_a^{u,t}&=F_{a,s}-g_{a,\mathrm{prev},I},&
 X_b^{u,t}&=F_{b,s}-g_{b,\mathrm{prev},I}.
 \end{aligned}                                                   \label{eq:hbase-X}
\end{equation}
Here \(g_{e,\mathrm{prev},I}\) is the last fit in slot \(e\) strictly
before \(t\), or zero.  In original slot order,
\begin{equation}
 H_I(F_{\cdot,s})
  =\sum_{\substack{1\le u<t\le n\\\iota(u)\ne\iota(t)}}
       H_I(X^{u,t}_I)+\mathrm{Rem}_{n,I},\qquad
 \sum_{I\in\mathcal U}|I||\mathrm{Rem}_{n,I}|\longrightarrow0     \label{eq:hbase-X-remainder}
\end{equation}
for every fixed finite \(\mathcal U\subseteq\mathcal I\).
The remainder is the term with four \(r_{e,n,I}\)'s plus, for each
slot \(e\), the term with \(g_{e,n,I}\) there and \(r_{l,n,I}\) in
the other three slots.

For sufficiently large \(A\), there is an absolute
\(0<\theta_0<A^{-7}<1\) such that, whenever \(t\ge12\), every pair
of factors \(X_l^{u,t}\) has form norm at most
\(e^{-2k_t^{\theta_0}}\).  All four factors have bounded normalized
\(L^2(I)\) size.
\end{lemma}

\begin{proof}
This is the algebra and per-cell comparison of the
\emph{The two-earliest-increments expansion} Lemma of
\cite[Section~5]{H}, now using the one operand \(F_{e,s}\) throughout
each cell.  For a finite \(n\), decompose each \(F_{e,s}\) into its
slot's increments \(g_t-g_{t-4}\) through \(n\) and the terminal
residual \(r_{e,n,I}\).  In a term with at least two increments,
record the earliest two levels \(u<t\).  In each other slot, the
sum of every later increment and its terminal residual is exactly
\(F_{e,s}-g_{e,\mathrm{prev},I}\).  Terms with at most one increment
sum to the stated remainder.  This proves the finite identity.

For updates \(v>w\) in distinct slots, replacing their pair of
projections by their actual inputs has form-norm error at most
\begin{equation}
 C_q\bigl(e^{-k_v}\exp(C_qk_w^{d_{\rm det}})+e^{-k_w}\bigr),     \label{eq:hbase-pair-comparison}
\end{equation}
where \(d_{\rm det}\) is absolute.  First test the residual at \(v\)
against the physical supremum bound of \(g_w\) as the \(L^q\)
opponent, then test the residual at \(w\) against the normalized
actual input.  Choose \(A>d_{\rm det}\) with fixed reserve.  The
same calculation with an actual input in one position makes the
form norm of two terminal residuals tend to zero.  The two remaining
factors of each remainder term have bounded \(L^2\) size.  Since
\(\mathcal U\) is finite, its absolute remainder sum tends to zero.

For the pair bound assume \(t\ge12\).  Each of \(X_j,X_a,X_b\)
has coefficient sum zero after its existing fitted operands are
replaced by the corresponding actual input: their preceding fits
occur at \(t-4\) and, for \(a,b\), among \(t-3,t-2,t-1\), so all exist.
Thus every pair contains a cancelling later factor; the early \(X_i\)
need not cancel when \(u\le4\).  If \(u<t-4\), consider a later factor
\(V\) paired with \(X_i=g_u-g_{u-4}\).  Each existing fitted operand
in \(V\) has index at least \(t-4\).  For each existing earlier fitted
operand \(g_w\), with
\(w\in\{u,u-4\}\), write the pairing with \(V\)
as the signed sum of pairings of \(g_w\) with the later residuals
\(g_v-F\); the actual-input terms cancel because those coefficients
sum to zero.  Testing these later residuals against \(g_w\) costs
at most \(C_qe^{-k_v}\exp(C_qk_w^{d_{\rm det}})\), and \(v>w\).
Pairs excluding \(i\) have no projection index below \(t-4\).

If \(u\ge t-4\), the distinct cyclic slots imply \(u\ge t-3\).
Every existing fitted operand in the tuple then has index at least \(t-7\).
For such a pair, and for the pairs excluding \(i\) in the preceding
case, expand into projection/actual pairs and apply
Equation~\eqref{eq:hbase-pair-comparison}.  The two-actual contribution
vanishes because at least one factor has coefficient sum zero.
The bounded number of errors is at most
\(C_qe^{-c_qk_{t-7}}\).  Since \(k_{t-7}=k_t^{A^{-7}}\), the stated
choice of \(\theta_0\) gives the required bound for large \(k_1\).
Contraction and the triangle inequality give the local sizes.
\end{proof}

\subsection{The high replacement target}

Fix one tuple from Equation~\eqref{eq:hbase-X}, put \(k=k_t\), and
write \(J_0\) for the normalized root on which its construction is
run (initially \(J_0=J\)).  Let \(\mathcal I_0\subseteq\mathcal I\)
be the original summands under consideration on \(J_0\).  The original
request-tree candidates are those already fixed for this construction.
Fix a common segment of its older fit families.  The positions \(i,j\) are now called low because their
two signed fitted differences will be expanded into original scalar
labels.  In the residual positions \(e=a,b\), we seek an integer
\(K_0\) and further inherited projections for which
\[
 X_e=F_{e,s}-g_{e,\mathrm{prev},I}
 \quad\hbox{can be replaced by}\quad
 P^{\rm new}_{e,s+K_0}F_{e,s}-g_{e,\mathrm{prev},I}
\]
with summable error.  A fixed label from each low position has a
common birth root: the deeper birth cell, or the segment root when
inherited there.  Their fitted coefficients remain outside the
scalar approximation.  The following tools are stated for a
generic fixed scalar pair, independently of the tuple notation.

\subsection{Two fixed scalar labels and their predictor list}

Fix distinct slots \(i,j\), let \(a,b\) be the complementary slots,
and fix bounded chart atoms \(\varphi_i,\varphi_j\) on their common
birth root \(J'\).  Their two-input form is
\[
 \mathcal L_I^{\varphi_i,\varphi_j}(z_a,z_b)
 :=|I|^{-2}\int_{\mathbb R^2}w_I(x,t)
       \varphi_i(y_i)\varphi_j(y_j)z_a(y_a)z_b(y_b)\,dx\,dt.
\]
For Hilbert-valued inputs, the displayed \(z_e\) denotes its physical
component.  The \emph{Progression separation} Lemma of
\cite[Section~4]{H} expands the scalar low product into terms
\(\psi_a(y_a)\psi_b(y_b)Q(t)\).  With the scalar coefficient removed,
\(\psi_a,\psi_b,Q\) are fixed bounded chart atoms of complexity at
most \(H\); unrestricted fixed linear modulations do not enter \(H\).
For one such term define
\[
 \begin{split}
 H_I(\psi_a z_a,\psi_b z_b;Q)
 :=|I|^{-2}\int_{\mathbb R^2}&w_I(x,t)Q(t)
       \psi_a(y_a)z_a(y_a)\\
       &\qquad\cdot\psi_b(y_b)z_b(y_b)\,dx\,dt .
 \end{split}
\]
Choose the cutoff allowed by the \emph{Single-scale separation} Lemma
of \cite[Section~4]{H} with
\[
 \zeta\in C_c^\infty(\mathbb R),\qquad
 0\le\zeta\le1,\qquad \zeta=1\ \hbox{on }[-1/3,1/3],
 \qquad \operatorname{supp}\zeta\subset(-1/2,1/2).
\]
This is permitted because \(x,x+3t\in I\) on the kernel support
imply \(|t|/|I|\le1/3\).  Set
\[
 A_r=\sup_{\xi\in\mathbb R}
       \left|\int_{\mathbb R}Q(ru)\zeta(u)e(-\xi u)\,du\right|.
\]
The normalization gives \(0\le A_r\le\|\zeta\|_1\le1\).
The plateau also gives \(\int\zeta\ge2/3\), so the source scalar
test retains a uniform constant.

For a dyadic \(K\subseteq J'\) and a selected finite use collection
\(\mathcal U\subseteq\mathcal I\) below \(K\), call an opposing
sequence \(z_{e,I}\) admissible on \(K\) if
\[
 z_{e,I}=\sum_{\ell=1}^{m_e}c_{e,\ell}V_{e,I}^{\ell},\qquad
 m_e\in\{0,1,2,3\},\qquad |c_{e,\ell}|\le1,
\]
where the scalar coefficients \(c_{e,\ell}\) are fixed independently
of \(I\) and \(s(I)\), and each \(V_{e,I}^{\ell}\) is one of these types:
\[
 F_{e,s(I)}|_I,\qquad P_{e,I}F_{e,s(I)},\qquad Z_{e,I}.
\]
All raw and projected constituents in slot \(e\) use the same full
normalized schedule \(F_{e,s}\); the two slots may have different
schedules.  Each projected constituent uses one inherited current
family common to all its cells, with the finitely many families under
a joint budget \(D_{\rm h}\).  Missing cells are completed by
inheritance and local refitting.  An energy constituent is supported
on its cell and satisfies the total bound
\[
 \sum_{I\in\mathcal U}|I|\|Z_{e,I}\|_{2,I}^2\le L^2|K|.
\]
Bilinearity and the triangle inequality reduce the estimates with two
opponents to at most nine constituent pairs, and the delayed estimate
to at most three opposing terms.  These absolute factors are absorbed
into constants.  All families used in an operator estimate are fixed
as realized families.

\begin{lemma}[A separated scale sum on depth inputs]
\label{lem:hbase-scale}
For the fixed separated term above, let both opponents be admissible
on \(K\) with parameters \(D_{\rm h},L\).  For
\(m\in\mathbb Z_{\ge0}\), let \(\mathcal U_m\subseteq\mathcal U\)
consist of any selected cells on which
\(e^{-m-1}<A_{|I|}\le e^{-m}\).  Then
\begin{equation}
 \sum_{I\in\mathcal U_m}|I|
   |H_I(\psi_a z_{a,I},\psi_b z_{b,I};Q)|
 \le e^{-m}\operatorname{poly}
       \bigl(1+m,\log(2H),\log(2D_{\rm h}),1+L\bigr)|K|.       \label{eq:hbase-scale-sum}
\end{equation}
These bands cover every positive \(A_{|I|}\).  If \(A_{|I|}=0\),
the corresponding two-input form vanishes.
\end{lemma}

\begin{proof}
Use the scalar construction and frequency estimates in the
\emph{A scale sum for one separated term} Lemma of
\cite[Section~5]{H}.  Put \(\mathcal K=1+m+\log(2H)\).
The \emph{Scale-wise Weyl alternative} and
\emph{Chart expansion and uniform tails} Lemmas of
\cite[Section~4]{H} produce phases
\[
 e(\lambda_*t^2+\gamma t),\qquad
 |\lambda_*|r^2\le D_1,\qquad \log D_1\le\operatorname{poly}(\mathcal K).
\]
The unscaled \(\lambda_*,\gamma\) are bounded-height rational
evaluations of one fixed underlying parameter vector.  This scalar
construction is independent of the opponents.

At one depth \(s\), disjointness, contraction, and
Equation~\eqref{eq:h-input-size} give
\[
 \sum_{\substack{I\in\mathcal U\\s(I)=s}}\|z_{e,I}\|_{L^2(I)}^2
 \le C(1+L)^2|K|.
\]
The energy case follows from its total square sum.  Cauchy--Schwarz
therefore bounds the sum of products of the two ordinary \(L^2\)
sizes by \(C(1+L)^2|K|\).  The original single-scale bilinear norm
is \(CA_r\), so this handles the polynomially many exceptional
scales in the source proof before expanding their coefficients.
It also handles the effective-quadratic and center-replacement
errors, whose scalar smooth-kernel weights have a small sum by the
\emph{Height bins} Lemma of \cite[Section~3]{H}.

For the chart tails, the pooled conclusion of the
\emph{Chart expansion and uniform tails} Lemma supplies a coefficient
envelope for each fixed pair \((\lambda_*,\gamma)\), uniformly over
the allowed rational splittings and bounded slow data at different
scales.  Split
\[
 e(\lambda_*t^2)=1+\bigl(e(\lambda_*t^2)-1\bigr).
\]
The second kernel has smooth-kernel norm at most
\(C|\lambda_*|r^2\operatorname{poly}(D_1)\).  For this one fixed
unscaled \(\lambda_*\),
\begin{equation}
 \sum_{\substack{r\ {\rm dyadic}\\|\lambda_*|r^2\le D_1}}
       |\lambda_*|r^2\le CD_1.                                \label{eq:hbase-quadratic-weight}
\end{equation}
The terms are zero when \(\lambda_*=0\); otherwise they form a
geometric series of ratio \(1/4\).  The per-depth estimate handles
this part.  As in the source scale-sum proof, choose the pooled mass
below \(e^{-2m}(1+D_1)^{-2}\), with fixed polynomial reserve for
the smooth-kernel costs, and then sum the fixed-pair envelopes.

For the linear parts we now instantiate the smooth-row Lemma.
In the coordinates \((y_a,y_b)\), the two one-variable marginals
of \(w_I\) vanish.  The double-primitive expansion in H's
\emph{A scale sum for one separated term} Lemma expresses it as a
summable series of tensors of interior, exactly mean-zero smooth
bumps.  Their fixed-order derivative bounds grow polynomially in
the full tensor index, while the tensor coefficients decay faster
than every fixed polynomial.

Fix one separated term, one opposing slot \(e\), and one full tensor
index.  For a retained frequency \(\nu\), the complete carrier is
\[
 m_\nu(y)=e(\nu y)\psi_e(y),\qquad \|m_\nu\|_\infty\le1,
 \qquad
 \mathcal W_{I,\nu}f
   =|I|^{-1/2}\int_I m_\nu(y)\omega_{I,\nu}(y)f(y)\,dy.
\]
The multiplier \(\psi_e\) is one fixed bounded function on the
common birth root.  Take the source-selected array on the original finite use
collection below \(J'\): it consists exactly of the retained intervals
and their retained clump representatives, with
at most one row for each \((I,\nu)\).  The source proof gives the
ordinary joint Bessel bound for this selected array, with squared
bound polynomial in \(\mathcal K\) and polynomial dependence on the
fixed tensor index.  Let \(B_H\) be its Bessel norm and
\(A_{\rm bump}\) its common value-plus-scaled-derivative bound.

The retained carrier catalog \(\mathcal N\) has
\(\log(2+\#\mathcal N)\le\operatorname{poly}(\mathcal K)\).
Indeed the chart dimension is at most \(\log H\); the logarithms of
the rational heights and pooled truncation bounds are polynomial
in \(\mathcal K\).  Counting the resulting bounded integer and
rational indices preserves a polynomial logarithm.  This is the
finite catalog supplied by the \emph{Chart expansion and uniform
tails} Lemma and the source pooled truncation.  Retain only rows
whose intervals lie in \(K\).  Zero extension of an \(L^2(K)\)
test to \(J'\) preserves their ordinary Bessel bound.  With
\(d_{\rm car}=2+\#\mathcal N\), Lemma~\ref{lem:smooth-depth} gives
\begin{equation}
 \left(\sum_{(I,\nu)\in\mathcal E}
      |\mathcal W_{I,\nu}f_{e,s(I)}|^2\right)^{1/2}
 \le C(B_H+A_{\rm bump})\log^2(2d_{\rm car})|K|^{1/2}.         \label{eq:h-smooth-use}
\end{equation}
Thus the raw scheduled input costs only the stated polynomial.
For the linear part of one tail pair, its frequency and multiplier
are fixed before its absolute coefficient envelope is summed;
the same argument uses its singleton complete carrier.  A fixed
predictor-list entry below is another singleton application.

The \emph{Compression} Lemma of \cite[Section~3]{H} is an ordinary
operator estimate.  If \(B_\nu=B_\nu\mathbf1_{u(\nu)}\) have joint
Bessel norm \(B\), and one common ordered list of inherited families
\(P^1,\ldots,P^m\), with \(m\ge1\), has joint budget \(D\), then
\begin{equation}
 \left(\sum_\nu
   \|B_\nu P^m_{u(\nu)}\cdots P^1_{u(\nu)}g\|^2\right)^{1/2}
 \le B\,\mathfrak C(m,D)\|g\|_2,\qquad
 \mathfrak C(m,D)=
 \bigl(C(m+1)^C\log(2D)\bigr)^{C(m+1)}.                       \label{eq:h-compression}
\end{equation}
This is uniform in repeated decision cells and labels on disjoint
branches.  A nonempty product has adjoint range in catalog
combinations on its whole decision cell; an empty product uses the
original Bessel bound.  Complements are expanded
into fixed ordered products, and averages are restored with their
total variation after compatible choices have been fixed.

For a projected constituent, fix this selected array and its one
inherited current family \(P_e\), together with its persistent
catalog of budget \(D_{\rm h}\).  Regard each physical row as a row
on the augmented space by ignoring private coordinates.  Its Bessel
bound is unchanged.  Compression with decision cell \(I\) gives the
ordinary norm \(B_H\mathfrak C(1,D_{\rm h})\) for
\(A_{I,\nu}=\mathcal W_{I,\nu}P_{e,I}\).  Moreover
\[
 A_{I,\nu}=A_{I,\nu}\mathbf1_I,\qquad
 A_{I,\nu}^*=P_{e,I}\mathcal W_{I,\nu}^*,
\]
whose output is a constant-coefficient combination of available
catalog columns on all of \(I\).  If the selected array is nonempty,
list its support depths increasingly as \(s_1<\cdots<s_N\), group
exactly its rows at each depth, and take both support and observation
partitions to be the full depth-\(s_n\) partition of \(K\).
Unused cells have no rows.  These partitions refine in order and
round \(n\) reads \(F_{e,s_n}\).  Lemma~\ref{lem:row-depth} therefore
gives
\begin{equation}
 \left(\sum_{(I,\nu)\in\mathcal E}
       |A_{I,\nu}F_{e,s(I)}|^2\right)^{1/2}
 \le C B_H\mathfrak C(1,D_{\rm h})\log(2D_{\rm h})|K|^{1/2}.   \label{eq:hbase-projected-rows}
\end{equation}
The empty array is trivial.  An energy constituent instead uses
H's per-interval Bessel bound over separated centers and its total
square sum.

The source Fourier-supremum test bounds each combined clump
coefficient by \(O(e^{-m})\); this scalar coefficient multiplies
the paired row arrays after their Bessel estimates.  Cauchy--Schwarz
between the two arrays and the rapid tensor coefficient decay prove
Equation~\eqref{eq:hbase-scale-sum}, including the linear pooled
tails treated with their separate singleton carriers.  If \(A_r=0\),
uniqueness of the Fourier transform gives \(Q(ru)\zeta(u)=0\)
almost everywhere.  Since \(\zeta=1\) on the relevant support,
the corresponding form is zero.
\end{proof}

\begin{corollary}[A finite predictor list for scheduled opponents]
\label{cor:hbase-predictor-list}
Fix the two bounded scalar atoms \(\varphi_i,\varphi_j\) on \(J'\),
of complexity at most \(H_{\rm low}\), and let \(B\ge1\).
There is a finite index set \(\mathfrak A\), a bound \(D_{\rm pred}\),
scalar coefficients \(c_{r,\alpha}\), and fixed triples
\((\psi_{a,\alpha},\psi_{b,\alpha},\lambda_\alpha)\) on \(J'\)
with \(\|\psi_{e,\alpha}\|_\infty\le1\), such that
\[
 \#\mathfrak A,\quad
 \operatorname{complexity}(\psi_{e,\alpha}),\quad
 |c_{r,\alpha}|\le D_{\rm pred},\qquad
 \log D_{\rm pred}\le
 \operatorname{poly}\bigl(1+B+\log(2H_{\rm low})\bigr).
\]
Here \(c_{r,\alpha}\) is a scalar depending only on the fixed pair,
its scalar construction, \(B\), and the dyadic scale \(r\); it is
zero unless \(|\lambda_\alpha|r^2\le D_{\rm pred}\).  The list and
these coefficients are chosen before any opposing sequences or
completed high catalogs.

Define its forms explicitly by
\begin{equation}
 \begin{split}
 T_I^\alpha(z_a,z_b)
 :=|I|^{-2}\int_{\mathbb R^2}&w_I(x,t)e(\lambda_\alpha t^2)
       \psi_{a,\alpha}(y_a)z_a(y_a)\\
       &\qquad\cdot\psi_{b,\alpha}(y_b)z_b(y_b)\,dx\,dt .
 \end{split}                                                   \label{eq:hbase-predictor-list}
\end{equation}
For every dyadic \(K\subseteq J'\), selected finite use collection
\(\mathcal U\) below \(K\), and pair of admissible opposing sequences
on \(K\), define the actual error \(e_I\) by the per-cell identity
\begin{equation}
 \mathcal L_I^{\varphi_i,\varphi_j}(z_{a,I},z_{b,I})
 =\sum_{\alpha\in\mathfrak A}c_{|I|,\alpha}
       T_I^\alpha(z_{a,I},z_{b,I})+e_I.                        \label{eq:hbase-list-identity}
\end{equation}
Then, uniformly over those opposing sequences and their completed
catalogs,
\begin{equation}
 \sum_{I\in\mathcal U}|I||e_I|
 \le e^{-B}\operatorname{poly}
       \bigl(\log(2D_{\rm h}),1+L\bigr)|K|.                    \label{eq:hbase-list-error}
\end{equation}
The error \(e_I\) may depend on the opposing sequence; the assertion
is its aggregate bound.
\end{corollary}

\begin{proof}
This is the \emph{A finite predictor list} Corollary of
\cite[Section~5]{H}, with Lemma~\ref{lem:hbase-scale} supplying its
opponent estimates.  The \emph{Progression separation} Lemma has
dyadic shell coefficient mass \(D_{\rm sep}T_{\rm sh}^{-100}\) and
shell complexity \((D_{\rm sep}T_{\rm sh})^C\), with
\(\log D_{\rm sep}\le\operatorname{poly}(\log(2H_{\rm low}))\).
The polynomial in Equation~\eqref{eq:hbase-scale-sum} is bounded by
a product of a polynomial in \(1+m+\log(2H)\) and one in
\(\log(2D_{\rm h}),1+L\).  The separation-shell cutoff and the
small-\(A_r\) cutoff may therefore be chosen using only
\(B,H_{\rm low}\).  The pooled-tail estimate supplies the required
absolute-sum precision on the retained bands.

The surviving bounded-height choices form the finite list after
distributing
\[
 e(\gamma t)
  =e\bigl(-\gamma y_a/(b-a)\bigr)
   e\bigl(\gamma y_b/(b-a)\bigr)
\]
into the fixed multipliers.  H's construction makes their scalar
coefficients depend only on scale and the fixed scalar data.
Normalize the multipliers by moving their scalar sizes to the
coefficients.  This gives the stated count, coefficient, and
complexity bounds and the identity with \(e_I\) its actual
difference.

The same fixed list has the stated bound on every \(K\subseteq J'\).
Use its original scalar cutoffs and carriers, and retain only the
source-selected rows whose cells are in the chosen use collection
inside \(K\).  Their ordinary Bessel bound follows by zero extending
a test on \(K\) to \(J'\), as in the preceding proof.  The restricted
schedules remain normalized, and the chosen current families are
locally refitted on \(K\).  Thus the per-depth, smooth-row, projected-row,
and energy estimates in that proof apply anew with root budget
\(|K|^{1/2}\) in each slot.  Summing the same discarded scalar
shells, bands, and pooled envelopes proves
Equation~\eqref{eq:hbase-list-error}.  This proves the aggregate
identity for all the stated opponents using one scalar list.
\end{proof}

\subsection{Delayed predictor replacement at a fixed use depth}

\begin{lemma}[Delayed predictor replacement on depth inputs]
\label{lem:hbase-delayed}
Use the fixed list and a dyadic \(K\subseteq J'\) from
Corollary~\ref{cor:hbase-predictor-list}.  Fix \(\alpha\in\mathfrak A\)
and a selected collection \(\mathcal U\) below \(K\) on which
\(|\lambda_\alpha||I|^2\le D_{\rm pred}\).  Let
\(K_0\in\mathbb Z_{\ge0}\) and \(0<\epsilon\le1\).
Suppose an inherited family \(P_{a,h}\) on the full partitions of
\(K\), with budget \(D_{\rm h}\), contains the same predictor
\[
 \Phi_v=(\overline{\psi_{a,\alpha}},\sqrt\epsilon\,\mathbf e_v)
\]
on every use cell and all its continued descendants.  Other columns
are allowed within that budget.  Let \(z_{b,I}\) be an admissible
opposing sequence on \(K\).  Then
\begin{align}
 &\sum_{I\in\mathcal U}|I|
 \left|T_I^\alpha(F_{a,s(I)},z_{b,I})
       -T_I^\alpha(P_{a,s(I)+K_0}F_{a,s(I)},z_{b,I})\right|
 \nonumber\\
 &\qquad\le
 C\bigl(2^{-K_0}+(1+K_0)\sqrt\epsilon\bigr)
 \operatorname{poly}\bigl(D_{\rm pred},\log(2D_{\rm h}),1+L\bigr)|K|.
                                                                    \label{eq:hbase-delayed}
\end{align}
For sequential replacement, let \(e\) denote the opposing slot and take
its full inherited family \(P_{e,h}\) on \(K\), under the stated joint budget.  On a use
cell \(I\) of depth \(s=s(I)\), the restriction of
\(P_{e,s+K_0}F_{e,s}\) to \(I\) is a current
constituent plus one whole rolling energy constituent with parameter
\(C(1+K_0)\log(2D_{\rm h})\).  Subtracting one older current
constituent still gives an admissible opponent.
\end{lemma}

\begin{proof}
We adapt the \emph{Delayed predictor replacement} Lemma of
\cite[Section~5]{H}.  Fix one use depth \(s\) throughout the
following calculation in \(h\), and abbreviate
\(\psi_{a,\alpha}\) by \(\psi_a\).  Let \(E_h\) be conditional
expectation on the full depth-\(h\) partition of \(K\), and set
\[
 r_h^{[s]}=\bigl(f_{a,s}-(P_{a,h}F_{a,s})_{\rm phys}\bigr)\psi_a,
 \qquad \mu_h^{[s]}=E_h r_h^{[s]}.
\]
Use these functions below a use cell, where \(\Phi_v\) is present.
If \(b_{h,v}^{[s]}\) is its coefficient on a depth-\(h\) cell,
orthogonality to \(\Phi_v\) and the zero private coordinates of
\(F_{a,s}\) give, with \(\Delta_h=P_{a,h+1}-P_{a,h}\),
\begin{equation}
 \mu_h^{[s]}=\epsilon b_{h,v}^{[s]},\qquad
 \mu_{h+1}^{[s]}-\mu_h^{[s]}
      =\sqrt\epsilon\,(\Delta_hF_{a,s})_v.                     \label{eq:hbase-predictor-moment}
\end{equation}
The second identity holds on the fine cells and survives every
further column insertion because the same predictor direction
persists.  Also
\[
 r_h^{[s]}-r_{h+1}^{[s]}
      =\psi_a(\Delta_hF_{a,s})_{\rm phys}.
\]

Let \(I\) have depth \(s\), put \(h_*=s+K_0\), and let \(\omega_I\)
be a smooth bump with bounded value and scale-normalized first
derivative.  For every finite \(n>h_*\),
\begin{align}
 \int_I r_{h_*}^{[s]}\omega_I
 ={}&\int_I\mu_{h_*}^{[s]}\omega_I
  +\sum_{h=h_*}^{n-1}\int_I
       (r_h^{[s]}-r_{h+1}^{[s]})(\omega_I-E_h\omega_I)
       \nonumber\\
 &+\sum_{h=h_*}^{n-1}\int_I
       (\mu_{h+1}^{[s]}-\mu_h^{[s]})
       (E_{h+1}\omega_I-E_h\omega_I)
  +\int_I r_n^{[s]}(\omega_I-E_n\omega_I).                    \label{eq:hbase-telescope}
\end{align}
Indeed, telescope
\(\int_I r_h^{[s]}(\omega_I-E_h\omega_I)\).  In the pairing with
\(E_{h+1}\omega_I-E_h\omega_I\), conditional expectation replaces
\(r_{h+1}^{[s]}\) by \(\mu_{h+1}^{[s]}\), while the pairing with
\(\mu_h^{[s]}\) vanishes.  The terminal term tends to zero because
\[
 \|r_n^{[s]}\|_{L^2(I)}\le2\|F_{a,s}\|_{L^2(I)},\qquad
 \|\omega_I-E_n\omega_I\|_{L^2(I)}
       \le C|I|^{1/2}2^{s-n}.
\]
For each finite \(n\), continue the family by retaining its labels
and locally refitting on the needed descendants.  These bounds then
dispose of the terminal term as \(n\) increases.

To sum the telescope across use depths, recall the following
projection energy bounds.  Let \(R\subseteq J\) be a normalized dyadic
subroot with its restricted normalized schedule \(F_{j,s}\).  Take one
full inherited continuation \((P_h)\) on its complete depth partitions
over the needed finite range, with catalog budget \(D\).  All depths
are the original depths in \(J\), with spatial origin \(s(R)\).  Put
\(\Delta_h=P_{h+1}-P_h\).  For a fixed integer offset \(\ell\), an
integer \(L_{\rm proj}\ge0\), and selected disjoint depth-\(s\) use cells
\(\mathcal U_s\) inside \(R\), Corollary~\ref{cor:row-projections} gives
\begin{equation}
 \begin{gathered}
 \sum_s\sum_{I\in\mathcal U_s}
   \|\mathbf1_I\Delta_{s+\ell}F_{j,s}\|_2^2
       \le C\log^2(2D)|R|,\\
 \sum_s\sum_{I\in\mathcal U_s}
   \|\mathbf1_I(P_{s+L_{\rm proj}}-P_s)F_{j,s}\|_2^2
       \le C(1+L_{\rm proj})^2\log^2(2D)|R|.
 \end{gathered}                                               \label{eq:h-projection-energy}
\end{equation}
Only valid projection depths are summed.  These are
Equations~\eqref{eq:row-increment} and \eqref{eq:row-delay}, and they
apply anew on every subroot.  The present Lemma uses \(R=K\) and
\(D=D_{\rm h}\).  In every finer evaluation belonging to a use at depth
\(s\), the operand remains \(F_{j,s}\).

We first sum the mean-zero tensor part.  Both conditional-expectation
errors in the two sums have supremum \(O(2^{s-h})\).  Normalize each
bump test by \(|I|^{-1/2}\); a length-weighted tensor form is a
product of two such tests, up to the fixed coordinate Jacobian.
At \(h=s+\ell\), the first normalized test is bounded by
\[
 C2^{-\ell}\|\mathbf1_I\Delta_{s+\ell}F_{a,s}\|_2,
\]
and the second has the additional factor \(\sqrt\epsilon\) from
Equation~\eqref{eq:hbase-predictor-moment}.  Apply the fixed-offset
part of Equation~\eqref{eq:h-projection-energy} at each \(\ell\)
and then Minkowski's inequality over \(\ell\ge K_0\).  The Hilbert
sum of these remainder tests is at most
\begin{equation}
 C2^{-K_0}\log(2D_{\rm h})|K|^{1/2}.                           \label{eq:hbase-predictor-remainder}
\end{equation}
This bound includes every fine cell of an increment and remains
valid after selecting the stated coarse uses.  The allowed bump
bounds multiply it.

When \(\int_I\omega_I=0\), the initial term may use
\(\mu_{h_*}^{[s]}-\mu_s^{[s]}\), since \(\mu_s^{[s]}\) is constant
on \(I\).  By Equation~\eqref{eq:hbase-predictor-moment}, it is a
sum of the \(K_0\) private increments at offsets \(0,\ldots,K_0-1\).
Equation~\eqref{eq:h-projection-energy} bounds the Hilbert sum of
these normalized tests by
\[
 CK_0\sqrt\epsilon\log(2D_{\rm h})|K|^{1/2}.
\]
Expand \(w_I\) into the mean-zero tensors used in the proof of
Lemma~\ref{lem:hbase-scale} and pair the error tests with the
opposing rows.  At one fixed tensor index, a raw opponent has the
singleton fixed carrier \(\psi_{b,\alpha}\) and one row per selected
use cell.  The singleton case of the ordinary mean-zero Bessel
calculation invoked in H's predictor proof, followed by
Lemma~\ref{lem:smooth-depth}, bounds it.  For a projected constituent,
apply the calculation in Equation~\eqref{eq:hbase-projected-rows}
to this fixed singleton array and its one inherited current family;
its full support and observation partitions and its use-depth reads
are the same.  An energy constituent uses the per-interval Bessel
bound and its total energy.  Rapid tensor coefficient decay absorbs
the polynomial bump costs.  This proves the asserted estimate for
the part of \(T_I^\alpha\) in which \(e(\lambda_\alpha t^2)\) is
replaced by one.

For the quadratic-difference part, hold \(s\) fixed.  Put
\(x_s=|\lambda_\alpha||I|^2\) at this depth.  A smooth tensor
expansion of \(w_I(e(\lambda_\alpha t^2)-1)\) has coefficient and
fixed-derivative cost at most
\[
 x_s\operatorname{poly}(D_{\rm pred}),\qquad x_s\le D_{\rm pred}.
\]
These bumps may have nonzero mean.  Ordinary orthogonality of the
full nested sequence on the one operand \(F_{a,s}\) gives
\[
 \sum_{h\ge h_*}\|\Delta_hF_{a,s}\|_2^2
 \le\|\mathbf1_KF_{a,s}\|_2^2.
\]
The local test bounds above and Cauchy--Schwarz in \(h\) consequently
give Hilbert-sum norm
\(C2^{-K_0}\|\mathbf1_KF_{a,s}\|_2\) for the two remainder sums
over the disjoint depth-\(s\) use cells.  For the initial mean,
\(\mu_{h_*}^{[s]}=\sqrt\epsilon(P_{a,h_*}F_{a,s})_v\);
contraction instead gives
\(C\sqrt\epsilon\|\mathbf1_KF_{a,s}\|_2\).
The opposing normalized tests at this depth have Hilbert-sum norm
at most \(C(1+L)|K|^{1/2}\), by disjointness, contraction, or the
energy bound.  Their total contribution is therefore at most
\[
 x_s\operatorname{poly}(D_{\rm pred})
       \bigl(2^{-K_0}+\sqrt\epsilon\bigr)(1+L)|K|.
\]
For this one fixed \(\lambda_\alpha\),
Equation~\eqref{eq:hbase-quadratic-weight} gives
\(\sum_{s:x_s\le D_{\rm pred}}x_s\le CD_{\rm pred}\).
This proves the quadratic-difference part.

Finally, on each use cell the delayed projection decomposes exactly as
\[
 P_{b,s+K_0}F_{b,s}
   =P_{b,s}F_{b,s}
     +(P_{b,s+K_0}-P_{b,s})F_{b,s}.
\]
The second term, restricted to each use cell as one whole constituent,
has total energy parameter at most \(C(1+K_0)\log(2D_{\rm h})\) by
Equation~\eqref{eq:h-projection-energy}.  Together with the current
first term it gives two admissible constituents; subtracting one older
current constituent gives the permitted third.
\end{proof}

\subsection{High fits and the replacement tuple}

Return to the fixed tuple and one common older segment
\(R_{\rm old}\subseteq J_0\).  We follow the
\emph{Truncating the two high positions} Proposition of
\cite[Section~5]{H}.  The low columns, their path counts, their
coefficient sums, and their atom complexities are bounded by
\(D_{\rm old}=\exp(k^{O(1)})\).  Expand the two low differences
only for this replacement.  Index the available scalar pairs by
\(p\), and let \(J_p\) be the deeper birth root, rebased to
\(R_{\rm old}\) when inherited there.  Pairs with disjoint births
have no common use.  Their coefficient products on a use cell
are bounded by a fixed power of \(D_{\rm old}\), and
\begin{equation}
 n_{\rm pair}(y):=\#\{p:y\in J_p\}\le D_{\rm old}^C,\qquad
 \sum_p|J_p|=\int_{R_{\rm old}}n_{\rm pair}(y)\,dy
       \le D_{\rm old}^C|R_{\rm old}|.                         \label{eq:hbase-low-packing}
\end{equation}

For every \(p\), choose its scalar list by
Corollary~\ref{cor:hbase-predictor-list}, with \(B\) a sufficiently
large fixed power of \(k\).  Choose one common numerical bound
\(D_{\rm pred}\) from that Corollary at this \(B\) and the common upper
bound on low atom complexity.  It bounds the count, coefficient sizes
at every scale, atom complexities, and quadratic cutoff of each
pair-specific list; its logarithm is a fixed polynomial in \(k\).
The lists remain distinct on their roots \(J_p\), and the predictor
count on a path uses only the pair roots on that path.  In high slot
\(e\), admit the conjugate predictors for each pair's list at that pair's
first joint availability, with distinct private directions of weight
\(\sqrt\epsilon\).  Pairs inherited at an attempt root have their
predictors present there; later pairs retain their predetermined
first-availability admissions.  The predictor path count is at most
\(D_{\rm old}^C D_{\rm pred}\).  After these lists and their common bound
are fixed, choose the integer \(K_0\) and \(\log(1/\epsilon)\) as
sufficiently large fixed powers of \(k\).

Here is the error rate that these choices must supply.  For a single
completed owner family with numerical budget \(D_{\rm h}\), put
\[
 \Lambda=C(1+K_0)\log(2D_{\rm h}),
\]
and
\[
 \begin{aligned}
 \mathfrak e_k(D_{\rm h})={}&
 e^{-B}\operatorname{poly}\bigl(\log(2D_{\rm h}),1+\Lambda\bigr)\\
 &+\bigl(2^{-K_0}+(1+K_0)\sqrt\epsilon\bigr)
       \operatorname{poly}
          \bigl(D_{\rm pred},\log(2D_{\rm h}),1+\Lambda\bigr).
 \end{aligned}
\]
The first term is the finite-list error on both sides of a change
of high input.  The second is the delayed error, summed over the
bounded list coefficients and entries.  In a sequential change,
the undelayed opponent is raw, optionally minus one previous current
constituent.
The delayed opponent is current plus one whole energy constituent with
parameter \(\Lambda\), optionally minus that one previous current
constituent.  Thus these are exactly the opposing classes of the two
scalar results.

Choose \(B\) to absorb the fixed low coefficient and pair-count
powers.  This fixes \(D_{\rm pred}\).  Choose \(K_0\) and
\(\log(1/\epsilon)\) afterward to absorb the fixed powers of
\(D_{\rm pred}\) and the same low costs, retaining an \(e^{-2k}\)
reserve.  Once the completed \(\log D_{\rm h}\) is a fixed power
of \(k\), its remaining factors in \(\mathfrak e_k\) are fixed
powers of \(k\).  Increasing the initial \(k_1\) then gives
\(D_{\rm old}^C\mathfrak e_k(D_{\rm h})\le Ce^{-k}\), for the fixed
power \(C\) needed below.  We now verify that the high catalogs
have this logarithmic size.

On one high attempt root \(R_{\rm fit}\), build
\(P^{\rm new}_{e,h}\) in projection-depth order.  The restricted
older data fix the previous approximants and the predictor admissions;
the high family's extra-fit catalog begins empty.  At each \(h\),
refine and locally refit the full preceding partition and admit all
predetermined predictors becoming available there.  This is done
also on branches with no further tests.  A candidate coarse use
cell at this step has depth \(s=h-K_0\); it is an assigned original
request-tree cell.  Depths with no such candidate still carry the
full refinements and admissions.

For a chosen precision \(\chi_e\), test on each candidate \(I\)
\[
 r_{e,I}=f_{e,s}-(P^{\rm new}_{e,h}F_{e,s})_{\rm phys}
\]
in \(\mathfrak d_{e,I}\) at threshold \(e^{-\chi_e}\).
The fine direct sum is an \(L^2(I)\) contraction, so this residual
has bounded normalized size.  Local detection gives on failure a
bounded atom \(\varphi\) with coarse moment at least \(\rho'\),
where \(\log(1/\rho')\) and its complexity log are bounded by a
fixed polynomial in \(k,\chi_e\).

Insert one common label \((\varphi,\mathbf e_v)\), born on \(I\),
into the mask of every depth-\(h\) subcell \(Q\subseteq I\), and
retain it on their descendants.  Its scalar function and private
direction are common, while its fitted coefficients may differ
between the \(Q\)'s.  With
\(m_Q=\fint_Q r_{e,I}\overline\varphi\), the fresh-column estimate
on each fine cell and Cauchy--Schwarz give
\begin{equation}
 c\sum_{\substack{Q\subseteq I\\s(Q)=h}}|Q||m_Q|^2
 \ge c|I|\left|\fint_I r_{e,I}\overline\varphi\right|^2
 \ge c(\rho')^2|I|.                                           \label{eq:hbase-high-gain}
\end{equation}
One request owns the distinct insertion rows on all these \(Q\)'s.
The label counts once at every point of \(I\).

After the admissions, perform simultaneous insertion rounds at \(h\).
Each still-failing coarse \(I\) may request one such label; every
other fine cell retains its projection.  The charged fine rows
have support and observation cells \(Q/Q\) in the full depth-\(h\)
partition and read \(F_{e,h-K_0}\).  Repeated rounds read the same
input.  The uncharged refinements and admissions remain in the
global nested sequence, and splitting on all fine cells preserves
Bessel norm one.  These data meet Lemma~\ref{lem:row-depth} and
the disjoint charged-row hypothesis of Lemma~\ref{lem:hbase-trial}.

The predetermined budget \(D_0\) includes all predictor admissions,
including future ones, and the bounded older data needed for the
joint budget.  Its logarithm is polynomial in \(k\).
A coarse path cap \(T\) adds at most \(T\) common bounded labels
of private weight one.  Thus \(D_T\le C D_0(T+1)^C\), and
\(\log(2D_T)\le k^{O(1)}+C\log(2T)\).
The one-trial Lemma chooses a cap with logarithm polynomial in
\(k,\chi_e\) and first-excess cells of total length at most
\(|R_{\rm fit}|/4\).  Only original coarse request cells are cut.

The two families at a cut have separate roles.

\emph{Old family.}  If a coarse cell at depth \(s\) is cut while
\(h=s+K_0\), every strict ancestor was already evaluated at a
smaller delayed depth.  On the cut branch, keep the old family's
current depth-\(h\) fine-cell masks unchanged through the remaining
rounds at \(h\).  On finer descendants inherit every label in those
masks, including extra-fit labels, and locally refit.  Its predictors
and their private directions persist, supplying every finer
projection in Equation~\eqref{eq:hbase-telescope} for the old
ancestor summands.  No further residual test on this branch is
needed for those summands.

\emph{New attempt.}  Send the cut coarse summands and their original
candidate descendants to a fresh attempt.  It retains the restricted
predetermined older data and the prescribed predictors: inherited
pairs have their predictors at the new root, and later pairs keep
their first-availability admissions.  Its extra-fit catalog has no
labels from the old extra fit.  This is the fresh-extra-catalog
choice permitted in the proof of the source high-truncation
Proposition.  The new \(D_0\) covers only its predetermined data
and admissions, and its cap counts only its own extra requests.

On each new root, restriction preserves the input normalization;
rebasing and local refitting preserve the predetermined path and
private budgets.  The same common-label gain, distinct fine rows,
and full-partition read chronology apply there, with the new empty
extra catalog.  Thus Lemma~\ref{lem:hbase-trial} applies anew and
again gives the quarter fraction.  The cut roots are proper original
request-tree descendants, and the finite recursion has total root
length at most \(4|R_{\rm fit}|/3\).  This realizes the continuation
and restart scheme of \cite[Section~5]{H} with the scheduled rows.

Apply this algorithm in the two high slots successively.  Let
\(S_{\rm old}=\exp(k^{O(1)})\) dominate the older supremum and
predictor costs.  Choose \(\chi_1\) with
\(e^{-\chi_1}S_{\rm old}\le e^{-2k}\), and construct the first high
fit and its capped catalog.  For a depth-\((s+K_0)\) cell \(Q\subset I\),
\[
 \|F_{e,s}\|_{2,Q}\le2^{K_0/2}\|F_{e,s}\|_{2,I}.
\]
The private-coordinate coefficient bound therefore bounds the first
fit's physical supremum on \(I\) by \(2^{K_0}\) times a fixed power
of its count and inverse private weights.  Its logarithm is
polynomial in \(k,\chi_1\); denote the resulting bound by \(S_1\).
Now choose \(\chi_2\ge\chi_1\) and construct the second fit so that
\[
 e^{-\chi_1}S_{\rm old}\le e^{-2k},\qquad
 e^{-\chi_2}S_1\le e^{-2k}.
\]
This also gives \(e^{-\chi_2}S_{\rm old}\le
e^{-\chi_1}S_{\rm old}\le e^{-2k}\).  Since \(\chi_1\) and \(\log S_1\)
are fixed powers of \(k\), both precision logs and both completed
catalog logs remain fixed powers of \(k\).  In particular their individual budgets have a common
numerical upper bound \(D_{\rm h}=\exp(k^{O(1)})\).  For any one
common owner piece below, a fixed enlargement of this number also
bounds the joint catalog of its two high families and the bounded
number of older families.  These
tuple-specific precision logs do not enter \(k_{t+1}=k_t^A\).

Intersect the two high ownership trees with the bounded number of
older segment trees.  Their maximal common refinement changes
owner only at a constituent segment root.  Its piece roots
\(\mathcal S(R_{\rm old})\) have total length at most
\(C|R_{\rm old}|\).  Each original summand in this older segment
is assigned to one piece \(S\), even when piece roots are nested.
On its assigned cells define
\begin{equation}
 Y_i=X_i,\qquad Y_j=X_j,\qquad
 Y_e=P^{\rm new}_{e,s+K_0}F_{e,s}
          -g_{e,\mathrm{prev},I}\quad(e=a,b).                  \label{eq:hbase-Y}
\end{equation}
Each projection here is the family of that cell's owner, with its
own full continuation.

We verify the aggregate replacement across these owners.  For a
low pair \(p\) and piece \(S\), put \(K=J_p\cap S\) when this is
nonempty; dyadic nesting makes \(K\) one cell.  Select only the
original coarse summands owned by \(S\) on which that pair is
available.  Use the original scalar list and coefficients chosen
on \(J_p\), and the one owner family in each high slot on \(S\),
with its complete descendant continuation.  On \(K\), rebase an
inherited pair and its predictors when needed and locally refit
the owner family.  This rebasing uses the same masks and retained spans
on each used cell, hence the same local owner projection there.
The scalar results apply anew on \(K\):
their source-selected raw rows are the original array restricted
to these cells, with the same Bessel bound by zero extension of
tests from \(K\), and their projected rows use this owner's
individual catalog.  The normalized schedules and row energies
have root budget \(|K|^{1/2}\).  Apply the finite-list identity
before and after each high change and the delayed estimate to the
retained entries.  With the low coefficient products outside,
the resulting cost for this pair and owner is at most a fixed
power of \(D_{\rm old}\) times
\(\mathfrak e_k(D_{\rm h})|K|\).

Writing \(n_{\rm seg}(y)=\#\{S\in\mathcal S(R_{\rm old}):y\in S\}\),
the intersection lengths satisfy
\[
 \sum_{p,S}|J_p\cap S|
 =\int_{R_{\rm old}}n_{\rm pair}(y)n_{\rm seg}(y)\,dy
 \le D_{\rm old}^C\sum_{S\in\mathcal S(R_{\rm old})}|S|.
\]
Here empty intersections contribute zero.  Combining the coefficient
and pair-count powers into \(D_{\rm old}^C\) proves
\begin{align}
 &\sum_{S\in\mathcal S(R_{\rm old})}
     \sum_{\substack{I\ {\rm owned}\\{\rm by}\ S}}
       |I|\,|H_I(X_I)-H_I(Y_I)|
 \nonumber\\
 &\qquad\le
 D_{\rm old}^C\mathfrak e_k(D_{\rm h})
       \sum_{S\in\mathcal S(R_{\rm old})}|S|
 \le Ce^{-k}|R_{\rm old}|.                                    \label{eq:hbase-high-errors}
\end{align}
This sums fresh estimates using each owner's continuation and the
original scalar list.

The high residuals pass their \(\mathfrak d\) tests on the coarse
cells.  The pair comparison used for \(X\) applies with the two high
precision logs ordered after the old logs.  If the earliest low
index is separated from \(t\), test the later residual difference
against each existing early fitted operand; otherwise every existing
old fitted operand has index at least \(t-7\).  For \(t\ge12\), each of
\(Y_j,Y_a,Y_b\) cancels when its existing fitted operands are replaced
by the same \(F_{e,s}\), so every pair has a cancelling later factor,
as in the proof of Lemma~\ref{lem:hbase-X-output}.  The high-truncation comparison in
\cite[Section~5]{H} therefore supplies an absolute exponent
\(\theta_{\rm av}>0\) for which each pair of \(Y\) factors has form
norm at most \(e^{-k^{\theta_{\rm av}}}\), when \(t\ge12\).
Contraction gives a fixed bound \(C_{\rm size}\) for every local
\(L^2(I)\) size.  Choose once
\begin{equation}
 0<\theta<\min\{1,\theta_{\rm av},\theta_0\},\qquad
 \eta=e^{-k^\theta}.                                          \label{eq:hbase-quality}
\end{equation}
After increasing \(k_1\), both the pair norm and
\(C_{\rm size}^2e^{-k^{\theta_{\rm av}}}\) are at most \(\eta\).
Thus every full \(Y\) pair has norm at most \(\eta\), and
\(|H_I(Y_I)|\le\eta\) for \(t\ge12\).

Across all common older segments, let \(\mathcal R_{\rm base}\) be the
final base segments and let \(\mathcal I_R\) be the preimage of \(R\)
under their owner map.  The sets \(\mathcal I_R\) partition this fixed
\(\mathcal I_0\), each
\(I\in\mathcal I_R\) lies in \(R\), and no auxiliary projection
cell is a summand.  Summing Equation~\eqref{eq:hbase-high-errors}
over the older segments gives the separate global statement
\begin{equation}
 \sum_{R\in\mathcal R_{\rm base}}|R|\le C|J_0|,\qquad
 \sum_{R\in\mathcal R_{\rm base}}\sum_{I\in\mathcal I_R}
       |I|\,|H_I(X_I)-H_I(Y_I)|\le Ce^{-k}|J_0|.              \label{eq:hbase-high-replacement}
\end{equation}

\subsection{The constructed tuple and its active patterns}

\begin{proposition}[Base data and active reduction]
\label{prop:hbase-data}
Fix the constructed tuple with later accuracy \(t\), put \(k=k_t\),
and use its final segments \(R\in\mathcal R_{\rm base}\) and original owner
sets \(\mathcal I_R\) from Equation~\eqref{eq:hbase-high-replacement}.
On an owned cell of depth \(s\), define
\begin{equation}
 \begin{gathered}
 Y_l=B_l+U_l,\qquad B_i=X_i,\quad B_j=X_j,\quad U_i=U_j=0,\\
 B_e=P^{\rm new}_{e,s}F_{e,s}-g_{e,\mathrm{prev},I},\qquad
 U_e=(P^{\rm new}_{e,s+K_0}-P^{\rm new}_{e,s})F_{e,s}
       \quad(e=a,b),\\
 \|B_{l,I}\|_{2,I},\ \|U_{l,I}\|_{2,I},\ \|Y_{l,I}\|_{2,I}\le C,\\
 \Lambda_k=C(1+K_0)\log(2D_{\rm base})\le k^C,\qquad
 \sum_{I\in\mathcal I_R}|I|\|U_{e,I}\|_{2,I}^2
       \le\Lambda_k^2|R|\quad(e=a,b).
 \end{gathered}                                               \label{eq:h-base-data}
\end{equation}
Each \(B_l\) is a signed sum of at most two current nested-catalog
projections applied to \(F_{l,s}\).  Their basic catalogs have
distinct private directions and inherited masks.  Their path counts,
column sizes, inverse private weights, bounded scalar atom
complexities, and current coefficient sums are bounded by
\(D_{\rm base}\), with
\[
 e^k\le D_{\rm base}\le\exp(k^C).
\]
For \(t\ge12\), the associated full \(Y\) tuple has the pair bound
and pointwise form bound \(\eta\) of Equation~\eqref{eq:hbase-quality}.

For three available basic scalar columns \(\boldsymbol\varphi\) in
distinct slots, let \(\operatorname{act}_I(\boldsymbol\varphi)\)
be their form norm against a normalized \(L^2(I)\) input in the
fourth slot.  For a sufficiently large fixed \(C_{\rm act}\), put
\[
 \mathcal A_R=
 \left\{I\in\mathcal I_R:
   \operatorname{act}_I(\boldsymbol\varphi)>D_{\rm base}^{-C_{\rm act}}
   \text{ for some available triple }\boldsymbol\varphi\right\}.
\]
This common active set has path multiplicity
\[
 \sum_{I\in\mathcal A_R}\mathbf1_I(y)\le M\le D_{\rm base}^C.
\]
Define three patterns in their original slots by
\[
 \begin{array}{lll}
 z^\varnothing_l=B_l&\text{for every }l,\\
 z^a_a=U_a,&z^a_l=B_l& (l\ne a),\\
 z^b_b=U_b,&z^b_l=B_l& (l\ne b),
 \end{array}
 \qquad \mathcal P=\{z^\varnothing,z^a,z^b\}.
\]
For \(z\in\mathcal P\), let
\begin{equation}
 \mathfrak M_R(z)
 :=\sum_{I\in\mathcal A_R}|I|
       \min\{\eta,|H_I(z_I)|\}.                               \label{eq:h-active-pattern-target}
\end{equation}
Then for \(t\ge12\),
\begin{equation}
 \sum_{I\in\mathcal I_R}|I|\,|H_I(Y_I)|
 \le C\eta\Lambda_k^2|R|+e^{-k}|R|
       +\sum_{z\in\mathcal P}\mathfrak M_R(z).                 \label{eq:h-active-minimum}
\end{equation}
The \(e^{-k}|R|\) term is the total inactive error for the three
patterns.  The finitely many leading accuracy levels have a bound
independent of the number of spatial depths.  These base data and
their estimates hold anew on normalized subroots with the inherited
basic lists and their local continuations.
\end{proposition}

\begin{proof}
The definitions give \(Y=B+U\).  Contraction bounds the local sizes.
Choose \(D_{\rm base}\) to dominate the bounded number of basic
catalogs used by this tuple and enlarge it to be at least \(e^k\).
The preceding logarithmic bounds give \(D_{\rm base}\le\exp(k^C)\).
On \(R\), use its owner families with their full inherited
continuations.  Equation~\eqref{eq:h-projection-energy}, applied
anew on this root, gives the stated \(U\) energies.

The scalar input for the active set is the
\emph{Few active scales for a fixed triple} Lemma of
\cite[Section~5]{H}.  For three fixed bounded basic scalar columns
on their common birth subtree and \(0<\delta\le1\), it gives on
every spatial path
\[
 \#\{I:\operatorname{act}_I(\boldsymbol\varphi)>\delta\}
 \le\operatorname{poly}\bigl(\log D_{\rm base}+\log(1/\delta)\bigr),
\]
and
\[
 \sum_{\substack{I\ {\rm on\ the\ path}\\
       \operatorname{act}_I(\boldsymbol\varphi)\le\delta}}
       \operatorname{act}_I(\boldsymbol\varphi)
 \le\delta\operatorname{poly}
       \bigl(\log D_{\rm base},\log(1/\delta)\bigr).
\]
The bounds are uniform in the path location and unrestricted real
oscillation parameters.  Every detected atom and predictor used
here is a fixed scalar label at birth.  At most \(D_{\rm base}^3\)
available triples lie on one path.  With
\(\delta=D_{\rm base}^{-C_{\rm act}}\), the first source bound gives
\(M\le D_{\rm base}^C\) after enlarging the fixed power.

On an inactive owned cell, expand three \(B\) factors of a pattern
into their basic scalar columns.  Their coefficients are constant
on that cell and their sums are bounded by a fixed power of
\(D_{\rm base}\); the fourth factor has bounded normalized \(L^2\)
size.  For each fixed triple, the second source bound controls the
inactive path sum by
\(D_{\rm base}^{-C_{\rm act}}
 \operatorname{poly}(\log D_{\rm base},C_{\rm act})\).
Bound the varying coefficients by their uniform sizes, sum the
pathwise triples and the three patterns, and integrate over path
points in \(R\).  The total absolute length-weighted inactive sum
is at most
\[
 D_{\rm base}^{C-C_{\rm act}}
       \operatorname{poly}(\log D_{\rm base},C_{\rm act})|R|
 \le e^{-k}|R|
\]
when \(C_{\rm act}\) dominates the fixed power costs and \(k_1\)
is sufficiently large.

It remains to derive the displayed minimum.  Let \(z^{UU}\) have
\(Y_i=B_i\) and \(Y_j=B_j\) in the low slots and \(U_a,U_b\) in
the other two.  For \(t\ge12\), the complementary \(Y_i,Y_j\)
pair bound and Cauchy--Schwarz give
\[
 \sum_{I\in\mathcal I_R}|I|\,|H_I(z_I^{UU})|
 \le C\eta
   \left(\sum_{I\in\mathcal I_R}|I|\|U_{a,I}\|_{2,I}^2\right)^{1/2}
   \left(\sum_{I\in\mathcal I_R}|I|\|U_{b,I}\|_{2,I}^2\right)^{1/2}
 \le C\eta\Lambda_k^2|R|.
\]
At each owned cell, multilinearity gives
\(H_I(Y)=\sum_{z\in\mathcal P}H_I(z)+H_I(z^{UU})\).
The full-tuple bound \(|H_I(Y)|\le\eta\), established before the
pattern split, therefore gives
\[
 |H_I(Y)|
 \le\min\left\{\eta,\sum_{z\in\mathcal P}|H_I(z)|
                         +|H_I(z^{UU})|\right\}
 \le |H_I(z^{UU})|
      +\min\left\{\eta,\sum_{z\in\mathcal P}|H_I(z)|\right\}.
\]
Sum over \(\mathcal I_R\), use the absolute inactive bound outside
\(\mathcal A_R\), and on the active cells apply
\(\min(\eta,\sum_z a_z)\le\sum_z\min(\eta,a_z)\).
This proves Equation~\eqref{eq:h-active-minimum}.  It assigns no
pair bound or pointwise \(\eta\) bound to an individual pattern.

For a leading level, the Local absolute estimate bounds the
two-\(U\) term by \(C\Lambda_k^2|R|\) without \(\eta\).
At each of the finitely many such levels, \(D_{\rm base}\),
\(\Lambda_k\), and \(M\) are constants independent of spatial
depth.  The local sizes and
\(\sum_{I\in\mathcal A_R}|I|\le M|R|\) bound the active contribution
independently of the number of depths.

Finally, on a normalized subroot the original schedule restrictions
remain normalized, basic masks remain inherited with the same
budgets, and local refitting gives the required full continuations.
The row estimate applied there supplies the \(U\) energies anew.
The arguments handed to localization are the whole \(B,U\) factors
above, together with the associated \(Y\) pair bound; the scalar
pair expansion was used only in the high replacement and the triple
expansion only on inactive cells.
\end{proof}
